%%
%% MapwiseFox -- SoftwareX Original Software Publication
%% Based on the official SoftwareX LaTeX template (elsarticle, numbered style).
%%

\documentclass[preprint,12pt, a4paper]{elsarticle}

%% Use the option review to obtain double line spacing
%% \documentclass[preprint,review,12pt]{elsarticle}

\usepackage{amssymb}
\usepackage{hyperref}
\setlength{\parindent}{0pt}

%% The lineno package adds line numbers; uncomment for review.
%\usepackage{lineno}

\journal{SoftwareX}

\begin{document}
\renewcommand{\labelenumii}{\arabic{enumi}.\arabic{enumii}}

\begin{frontmatter}

%% ---------------------------------------------------------------------
%% Title: "SoftwareName: Then a short title" per the template convention.
\title{MapwiseFox: A Modular Toolkit for Systematic Literature Reviews}

%% AUTHORS: given name(s) and family name(s), order matching the submission
%% system. Use \corref for the corresponding author.
%% MISSING: confirm the full author list and order before submission.
\author{Andrei Olar}

%% AFFILIATIONS: full name and postal address, including country.
%% MISSING: confirm the affiliation address.
\address{[Affiliation, Department, Street, City, Postal code, Country]}

%% CORRESPONDING AUTHOR: keep the email address current through publication.
%% MISSING: add the corresponding author email.
\ead{[corresponding.author@example.org]}

%% ---------------------------------------------------------------------
%% ABSTRACT: ca. 100 words, concise and factual, no references.
%% Must state purpose, principal results, and major conclusions.
\begin{abstract}
MapwiseFox is an open-source toolkit that supports the main stages of
systematic literature reviews and systematic mapping studies. It runs a
single search query across six bibliographic databases, deduplicates the
results, performs backward and forward snowballing, splits workloads
among reviewers, computes inter-rater agreement metrics, and offers
large-language-model-assisted study selection and quality assessment. A
web application lets reviewers manually select primary studies against
customisable inclusion and exclusion criteria. The toolkit is
self-hostable, ships as a single Docker image, and is released under the
MIT licence.
\end{abstract}

%% KEYWORDS: 1 to 7, English, avoid multi-word phrases with "and"/"of".
\begin{keyword}
Systematic literature review \sep
Systematic mapping study \sep
Research synthesis \sep
Information retrieval \sep
Snowballing \sep
Large language model tools
\end{keyword}

\end{frontmatter}

%\linenumbers

%% =====================================================================
%% Required Metadata
%% =====================================================================

\section*{Current code version}

%% Ancillary data table. Replace the right column with the correct
%% information; leave the left column as in the template.

\begin{table}[!h]
\begin{tabular}{|l|p{6.5cm}|p{6.5cm}|}
\hline
\textbf{Nr.} & \textbf{Code metadata description} & \textbf{Please fill in this column} \\
\hline
C1 & Current code version & v0.8.1 \\
\hline
C2 & Permanent GitHub link to code/repository used for this code version & \url{https://github.com/koosie0507/mapwisefox} \\
\hline
C3 & Legal Code License & MIT \\
\hline
C4 & Code versioning system used & git \\
\hline
C5 & Software code languages, tools, and services used & Python, TypeScript, React, FastAPI \\
\hline
C6 & Compilation requirements, operating environments \& dependencies & Python 3.13, uv, Node.js, Docker \\
\hline
C7 & If available Link to developer documentation/manual & %% MISSING: add docs URL if hosted
\\
\hline
C8 & Support email for questions & %% MISSING: add support email
\\
\hline
\end{tabular}
\caption{Code metadata (mandatory)}
\label{tab:metadata}
\end{table}

%% =====================================================================
%% Main text
%% =====================================================================
%% Maximum 6 pages (excluding metadata, tables, figures, references)
%% with a 4000-word limit. Five mandatory sections follow.

%% ---------------------------------------------------------------------
%% 1. Motivation and significance
%% ---------------------------------------------------------------------
%% Explain why the software is important, what problem(s) it solves,
%% how it contributes to scientific discovery, how a user runs it, and
%% what related work exists.

\section{Motivation and significance}

Systematic literature reviews (SLRs) and systematic mapping studies
(SMSs) are evidence-synthesis methods used across many disciplines,
from software engineering to medicine. They follow a defined protocol:
search multiple databases, remove duplicates, expand the seed set
through snowballing, screen studies against inclusion and exclusion
criteria, assess quality, and extract data \cite{kitchenham2007,petersen2008,page2021}.
Each stage is labour-intensive, and the volume of records can reach
thousands.

%% [Problem] The exact problem MapwiseFox solves.
Practitioners typically rely on a mix of commercial platforms and
ad-hoc scripts, moving data between formats manually. This is
error-prone, hard to reproduce, and limits the portability of review
protocols across teams and institutions.

%% [Contribution] How the software contributes to scientific discovery.
%% MISSING: cite a research paper that uses MapwiseFox once available.
MapwiseFox addresses this gap by providing an integrated, open-source
toolkit that covers the main stages of an SLR or SMS up to data
extraction. It unifies the pipeline under a single data model and ships
as a self-contained Docker image, making review protocols portable and
reproducible.

%% [Experimental setting] How a user uses the software.
Users interact with MapwiseFox through command-line tools for search,
deduplication, snowballing, workload splitting, and metrics, and
through a web application for manual study selection. A typical
workflow starts with a search across all configured backends, followed
by deduplication, optional snowballing, workload splitting among
reviewers, and finally selection via the web interface or the
LLM-assisted assistant.

%% [Related work] Introduce related work in literature.
%% MISSING: add citations to comparable tools and methodological
%% guidelines (e.g., ASReview [1], Rayyan, Covidence, Publish or Perish,
%% PRISMA [2], Kitchenham [3], Petersen [4]).
Related work includes both commercial platforms (Covidence,
Rayyan~\cite{rayyan}) and open-source tools
(ASReview~\cite{van2021}). MapwiseFox differs by covering the full
pipeline from search to selection in a single self-hostable package,
and by integrating large-language-model-assisted selection alongside
traditional manual review.

%% ---------------------------------------------------------------------
%% 2. Software description
%% ---------------------------------------------------------------------
%% Describe the software with enough detail for the reader to understand
%% its impact.

\section{Software description}

%% 2.1 Software architecture
\subsection{Software architecture}

MapwiseFox is organised as a Python 3.13 monorepo managed by a single
\texttt{uv}~\cite{uv} workspace. Eight packages share the \texttt{mapwisefox}
namespace: \texttt{search}, \texttt{deduplication}, \texttt{snowballing},
\texttt{split}, \texttt{metrics}, \texttt{assistant},
\texttt{common-config}, and \texttt{web/backend}. The web frontend is a
standalone React and Vite application outside the Python workspace.

All packages operate on Microsoft Excel files as the interchange format,
except \texttt{search} and \texttt{web}, which support additional
formats. A shared configuration package keeps behaviour consistent
across tools. A Dockerfile builds a single image containing every CLI
and the web application. In production, Caddy serves the
React~\cite{react} single-page application and reverse-proxies API
and authentication routes to the FastAPI~\cite{fastapi} backend.

%% Architecture overview (Mermaid source; to be rendered as Fig. 1).
%% The diagram shows the top-level packages, the shared data model,
%% external services, and the flow of records through the pipeline.
%%
%% %%{init: {'theme': 'base', 'themeVariables': { 'primaryColor': '#e6f2ff', 'primaryTextColor': '#0d1b2a', 'primaryBorderColor': '#1b4965', 'lineColor': '#1b4965', 'secondaryColor': '#fff4e6', 'tertiaryColor': '#e8f5e9'}}}%%
%% flowchart TB
%%     subgraph External["External systems"]
%%         WOS[Web of Science]
%%         Scopus[Scopus]
%%         SD[ScienceDirect]
%%         Springer[Springer]
%%         ACM[ACM Digital Library]
%%         IEEE[IEEE Xplore]
%%         S2[Semantic Scholar API]
%%         LLM[OpenAI-compatible LLM endpoints]
%%     end
%%
%%     subgraph CLI["CLI packages (uv workspace)"]
%%         SearchPkg[search]
%%         DedupPkg[deduplication]
%%         SnowPkg[snowballing]
%%         SplitPkg[split]
%%         MetricsPkg[metrics]
%%         AssistPkg[assistant]
%%         ConfigPkg[common-config]
%%     end
%%
%%     subgraph WebStack["Web application"]
%%         Caddy[Caddy reverse proxy]
%%         React[React / Vite frontend]
%%         FastAPI[FastAPI backend<br/>web/backend]
%%         DB[(Persistent store)]
%%     end
%%
%%     subgraph Data["Shared data model"]
%%         Excel[(Microsoft Excel workbooks)]
%%         Common[Common record schema]
%%     end
%%
%%     Query([Review protocol / query DSL]) --> SearchPkg
%%     SearchPkg -->|translates to native syntax| WOS & Scopus & SD & Springer & ACM & IEEE
%%     WOS & Scopus & SD & Springer & ACM & IEEE -->|raw exports| SearchPkg
%%     SearchPkg -->|normalised records| Excel
%%     Excel --> DedupPkg
%%     DedupPkg -->|deduplicated seed set| Excel
%%     Excel --> SnowPkg
%%     SnowPkg -->|DOI lookup| S2
%%     S2 -->|citations / references| SnowPkg
%%     SnowPkg -->|snowball candidates| Excel
%%     Excel --> SplitPkg
%%     SplitPkg -->|reviewer partitions| Excel
%%     Excel --> MetricsPkg
%%     MetricsPkg -->|agreement / IR metrics| Reports([Report files])
%%     Excel --> AssistPkg
%%     AssistPkg -->|rendered Jinja2 prompts| LLM
%%     LLM -->|include/exclude / quality scores| AssistPkg
%%     AssistPkg -->|annotated records| Excel
%%     Excel --> FastAPI
%%     FastAPI -->|REST API| React
%%     React -->|manual selection UI| Reviewer([Human reviewer])
%%     Reviewer -->|inclusion / exclusion decisions| React
%%     ConfigPkg -.->|shared settings| SearchPkg & DedupPkg & SnowPkg & SplitPkg & MetricsPkg & AssistPkg & FastAPI
%%     Caddy -->|static SPA| React
%%     Caddy -->|/api /auth| FastAPI
%%     FastAPI --> DB
%%
%% Figure 1: Top-level architecture of MapwiseFox. CLI packages (left)
%% operate on a shared Excel-based data model and are driven by a common
%% configuration package. The web application (right) exposes the same
%% data model through a FastAPI backend and React frontend. External
%% bibliographic databases feed the search and snowballing stages, while
%% OpenAI-compatible LLM endpoints power the assistant.

%% 2.2 Software functionalities
\subsection{Software functionalities}

%% [search] One query, six backends.
\textbf{Search.} Accepts a single query expressed in a small
domain-specific language and translates it into the native query syntax
of six backends: Web of Science, Scopus, ScienceDirect, Springer, ACM
Digital Library, and IEEE Xplore. Results are normalised into a common
record format and persisted via pluggable adapters.

%% [deduplication]
\textbf{Deduplication.} Clusters paper metadata coming from multiple
sources so that the same record appearing in several search exports is
collapsed into a single entry.

%% [snowballing]
\textbf{Snowballing.} Performs configurable-depth backward or forward
snowballing starting from a set of primary studies identified by DOI.
It resolves citations and references through the Semantic Scholar API
and writes unique new candidates to Excel workbooks.

%% [split]
\textbf{Split.} Partitions a review workload across multiple
participants using several schemes, so that each reviewer receives a
disjoint subset of the studies.

%% [metrics]
\textbf{Metrics.} Computes agreement and relevance metrics used
throughout the review process: Cohen's Kappa for categorical ratings,
intraclass correlation coefficient (ICC) and concordance correlation
coefficient (CCC) for continuous ratings, and information-retrieval
measures for search quality.

%% [assistant]
\textbf{Assistant.} Uses large language models to perform study
selection (include or exclude) or quality assessment based on
user-defined criteria. It supports any OpenAI-compatible provider
(OpenAI, Anthropic, Google, AWS Bedrock, and local endpoints) and
renders prompts from Jinja2 templates.

%% [web]
\textbf{Web.} A FastAPI backend and React frontend that let reviewers
manually select primary studies against customisable inclusion and
exclusion criteria. Authentication is optional and uses OpenID Connect
when enabled.

%% MISSING: the assistant currently supports only one LLM provider at a
%% time. Multi-provider routing is not yet implemented.
%% MISSING: there is no built-in data-extraction stage; the pipeline
%% ends at study selection and quality assessment.
%% MISSING: the web application does not yet support multi-reviewer
%% concurrency controls or conflict-resolution workflows.

%% 2.3 Sample code snippets analysis (optional)
\subsection{Sample code snippets}

%% A typical end-to-end run, citing the exact commands.
A review starts with a search across all configured backends, after
which the exported results are deduplicated, optionally snowballed, and
split among reviewers:

\begin{verbatim}
uv run search -D ./data/run-1
uv run deduplicate -i ./data/run-1 -o ./data/dedup.xlsx
uv run snowball -i ./data/dedup.xlsx -d backward --depth 2
uv run split-workload -i ./data/dedup.xlsx -n 3
uv run web --host 127.0.0.1 --port 8000
\end{verbatim}

%% MISSING: verify the exact CLI invocations against the current code
%% before submission, as flags may have changed.

%% ---------------------------------------------------------------------
%% 3. Illustrative examples
%% ---------------------------------------------------------------------
%% Provide at least one illustrative example to demonstrate the major
%% functions. OPTIONAL: an explanatory video or screencast (MP4, max
%% 150 MB, 640x480, 30 fps) may be submitted as supplementary material.

\section{Illustrative examples}

%% MISSING: provide a concrete walk-through using a real or sample
%% dataset. Show the input, the command, and the output for at least
%% the search and deduplication stages. Consider including a screenshot
%% of the web application.

%% ---------------------------------------------------------------------
%% 4. Impact
%% ---------------------------------------------------------------------
%% This is the main section; reviewers weight it appropriately.
%% Indicate new research questions enabled, improvements to existing
%% ones, changes to daily practice, adoption, and commercial use.
%% Points 1 and 2 are best demonstrated by citable publications.

\section{Impact}

%% [New research questions]
By integrating LLM-assisted selection with traditional manual review,
MapwiseFox enables research into human--AI collaboration in
evidence-synthesis workflows, including questions about agreement,
bias, and throughput when large language models assist with study
screening.

%% [Improving existing research questions]
The toolkit reduces the manual effort required at each stage of an
SLR or SMS. The search DSL lets a researcher write one query and run
it against six databases, eliminating the need to translate syntax
manually. The deduplication and snowballing stages further reduce
redundant work.

%% [Daily practice]
MapwiseFox changes the daily practice of its users by replacing a
chain of disconnected tools and manual data transfers with a single,
scriptable pipeline. The Docker image makes the toolkit available to
researchers who cannot install each dependency individually.

%% [Adoption]
%% MISSING: report download counts, user numbers, or citable
%% publications once the software has been publicly released.

%% [Commercial use]
%% MISSING: report any commercial use or spin-off companies if
%% applicable.

%% ---------------------------------------------------------------------
%% 5. Conclusions
%% ---------------------------------------------------------------------
%% Mandatory section summarising the software.

\section{Conclusions}

MapwiseFox provides an integrated, open-source toolkit that covers the
main stages of a systematic literature review or systematic mapping
study up to data extraction. It runs a single query across six
bibliographic databases, deduplicates results, performs
configurable-depth snowballing, splits workloads among reviewers,
computes agreement metrics, and offers both LLM-assisted and
web-based study selection. The toolkit is self-hostable, ships as a
single Docker image, and is released under the MIT licence.

%% MISSING: finalise the future-work outlook (data-extraction stage,
%% multi-provider LLM routing, conflict-resolution workflows, additional
%% search backends).

%% =====================================================================
%% Acknowledgements
%% =====================================================================
%% Optionally thank people and institutes. Place directly before the
%% reference list, not on the title page.
%% MISSING: add acknowledgements if applicable.

\section*{Acknowledgements}
[Acknowledgements text, if any.]

%% =====================================================================
%% References
%% =====================================================================
%% Numbered (Vancouver-style) references via BibTeX.
%% Compile with: pdflatex -> bibtex -> pdflatex -> pdflatex.

\bibliographystyle{elsarticle-num}
\bibliography{references}

%% =====================================================================
%% Current executable software version
%% =====================================================================
%% This section is required by the SoftwareX template.
%% MISSING: add a Zenodo DOI for the archived executable version once
%% the release is minted.

\section*{Current executable software version}
[To be completed with the Zenodo DOI and archival details upon release.]

\end{document}
